\documentclass[journal,10pt]{IEEEtran}

% Packages
\usepackage{amsmath,amssymb,amsfonts}
\usepackage{algorithmic}
\usepackage{graphicx}
\usepackage{textcomp}
\usepackage{xcolor}
\usepackage{booktabs}
\usepackage{hyperref}
\usepackage{cite}

\begin{document}

\title{Deforestation Detection using Deep Learning and Multispectral Sentinel-2 Satellite Imagery}

\author{Your Name, \IEEEmembership{Member, IEEE}
\thanks{Manuscript received June 27, 2026. Corresponding author: Your Name (email: author@university.edu).}}

\markboth{IEEE Transactions on Geoscience and Remote Sensing,~Vol.~XX, No.~X, June~2026}%
{Your Name \MakeLowercase{\textit{et al.}}: Deforestation Detection using Deep Learning on Sentinel-2 Imagery}

\maketitle

\begin{abstract}
Rapid deforestation poses severe threats to global ecosystems, biodiversity, and carbon sequestration cycles. In this study, we propose a scalable deep learning framework for precise forest cover segmentation and change detection. Leveraging multispectral Sentinel-2 satellite bands (RGB and Near-Infrared), we implement a modified convolutional segmentation network (U-Net) benchmarked against the EuroSAT land cover dataset. Furthermore, we employ a bi-temporal change detection architecture validated against Global Forest Watch (GFW) historical alerts to detect tree cover loss events. Our approach yields high Intersection-over-Union (IoU) scores, demonstrating the viability of automated satellite telemetry for real-time ecological monitoring.
\end{abstract}

\begin{IEEEkeywords}
Deforestation detection, deep learning, Sentinel-2, EuroSAT, Global Forest Watch, change detection, semantic segmentation.
\end{IEEEkeywords}

\section{Introduction}
\IEEEPARstart{F}{orests} cover approximately 31\% of the global land area, acting as critical carbon sinks and hosting major portions of terrestrial biodiversity. However, anthropogenic pressures, commercial logging, and agricultural expansion have accelerated canopy loss worldwide...

\section{Related Work}
\subsection{Land Cover Classification with EuroSAT}
Land cover classification benchmarks have historically progressed from coarse resolution sensors to high-resolution multispectral imagery. The EuroSAT dataset provides...

\subsection{Deep Learning for Satellite Imagery Segmentation}
Semantic segmentation models like U-Net, DeepLabv3+, and Feature Pyramid Networks (FPNs) have become standard for pixel-wise mapping of forest boundaries...

\section{Methodology}
Our end-to-end framework consists of three main components: (1) multi-spectral patch generation, (2) semantic forest cover segmentation, and (3) bi-temporal deforestation change detection.

\subsection{Multispectral Image Preprocessing}
We utilize the Blue (B02), Green (B03), Red (B04), and Near-Infrared (B08) bands of Sentinel-2 satellite imagery. High-resolution tiles are partitioned into local patches of dimension $P \times P$:
\begin{equation}
I \in \mathbb{R}^{P \times P \times 4}
\end{equation}

\subsection{Forest Segmentation Model}
A U-Net encoder-decoder network is trained to classify each pixel as either forest or non-forest...

\subsection{Bi-Temporal Change Detection}
To detect deforestation events, two co-registered image patches $I_{t1}$ and $I_{t2}$ taken at separate dates are passed through our change detection module...

\section{Experiments and Evaluation}
\subsection{Datasets and Labels}
- \textbf{EuroSAT:} Pre-training benchmark.
- \textbf{Sentinel-2:} Input satellite imagery source.
- \textbf{Global Forest Watch (GFW):} Deforestation validation masks.

\subsection{Evaluation Metrics}
To evaluate segmentation fidelity, we report the Intersection-over-Union (IoU) and F1-score (Dice coefficient):
\begin{equation}
\text{IoU} = \frac{\text{TP}}{\text{TP} + \text{FP} + \text{FN}}
\end{equation}

\section{Results and Discussion}
% Quantitative comparison tables and confusion matrices will go here

\section{Conclusion}
This paper presented an integrated machine learning workflow for remote sensing-based forest loss assessment...

\bibliographystyle{IEEEtran}
\bibliography{references}

\end{document}
